\chapter*{怎样读这一卷}
\addcontentsline{toc}{chapter}{怎样读这一卷}

本卷从隋受禅的五八一年写到宋代取代后周的九六〇年。隋、唐与五代十国相连而不相互吞没：唐亡以后，吴、南唐、吴越、闽、南汉、楚、荆南、前后蜀、北汉等与中原五朝并行；吐蕃、回鹘、渤海、新罗、日本和西域政权也按事件进入自己的位置。

先用年号、干支、改元和并行政权的日期建立次序，再看运河、道路、军镇、税源、城市与边疆如何给选择设限。把结构性的资源与制度条件、近处的触发、行动所依赖的组织和补给，以及直接和中期后果分别阅读；“盛”“衰”或“治”“乱”不能替代这种因果检验。

《隋书》、两《唐书》、两《五代史》与《资治通鉴》及《考异》构成编年脊柱；《通典》《唐会要》、诏令、敦煌吐鲁番文书、碑志、城址和外部材料补充制度与地方尺度。官修褒贬、制度条文、军数、朝贡语言和一地考古材料都有边界，不能直接推出全国执行、固定族群或精确疆界。

页中的 \texttt{event} 是首次深描的入口，\texttt{turningpoint} 提醒哪一步改变了后续路径，\texttt{causalchain} 拆解条件、触发与反馈；\texttt{textwindow} 交代引文的语境，\texttt{sourcenote}、\texttt{debate} 说明可证与未决之处。它们服务于连续阅读，而不是把多政权世界压成几条结论。
